% Rhetorical-Ontological Diachronic Analysis (RODA) -- Worked Walkthrough: the RDR2 saloon brawl.
% Compile with pdflatex (TeX Live >= 2018) or xelatex/lualatex. Greek is transliterated (italic, macrons).
% Sources: CANONICAL-METHODOLOGY-v5.md + VALIDATION-v5/v4/v3-barfight.md.
\documentclass[11pt]{article}
\usepackage[utf8]{inputenc}
\usepackage[T1]{fontenc}
\usepackage{lmodern}
\usepackage[margin=1in]{geometry}
\usepackage{enumitem}
\usepackage{titlesec}
\usepackage{xcolor}

\setlist[itemize]{leftmargin=1.5em,itemsep=3pt,topsep=3pt,parsep=0pt}
\setlist[itemize,2]{leftmargin=1.4em,itemsep=1.5pt,topsep=1.5pt,label=$\circ$}
\titlespacing{\section}{0pt}{13pt}{5pt}
\newcommand{\inb}{\textbf{Brawl:}~}
\newcommand{\too}{\ensuremath{\to}}

\title{\textbf{Rhetorical--Ontological Diachronic Analysis (RODA)}\\[5pt]
\large A Worked Walkthrough---the Saloon Brawl in \emph{Red Dead Redemption~2}}
\author{Dalton Salvo}
\date{June 2026}

\begin{document}
\maketitle

\noindent\textbf{What the method does.} RODA analyzes player behavior as the \emph{actualization of desire}: it traces how a designed world moves a disposed, desiring being from a perceptual encounter, through a charged mental image and a committed, emotionally-laden judgment, to a completed action---and how that action sediments into the player's lasting disposition (\textit{hexis}). One applies it by walking a bounded episode once through the chain $A_0\too A_4$, naming at each step where the designed world exerts pressure and where the player's own desire does the work, diagnosing the substantial union the world drives toward (\emph{rhetorical incorporation}), and reading the motive, cause, and dispositional residue. Its distinctive yield is an account that is at once ontological, ecological, phenomenological, and social: it explains not merely \emph{that} a player is involved but \emph{why} and \emph{how} a designed world remakes who they are---installing skills, cultivating or corrupting character, and, at the limit, altering how they dwell in the world beyond the game.

\medskip
\noindent\textbf{The episode.} Early in \emph{Red Dead Redemption~2} the player, as the outlaw Arthur Morgan, drinks with his gang in a Valentine saloon. A gang-mate, Bill, shoves a stranger, and a brawl erupts that spills into the muddy street, where a crowd gathers and eggs on a local fighter, Tommy. The game teaches combat in real time (an on-screen prompt, ``Hold R to block''); Arthur trades blows, breaks a grapple, and beats Tommy down---until a bystander begs him to stop (``I beg you''), opening a small fork of restraint that the Honor system will later compound. The encounter is brief, scripted to be won, and almost wholly a matter of belonging: Arthur fights not from a private grievance but because he is one with the gang.

\medskip
\noindent\textbf{Rhetorical incorporation.} The method's central concept names the substantial union a designed world drives the actualization toward: the player becomes one-with the world they act \emph{in} and one-with the avatar and others they act \emph{as} and \emph{with}, while remaining a distinct self. It is built by fusing and re-grounding three accounts. \emph{Calleja's incorporation} supplies the phenomenology of a player absorbed into, and embodied within, a gameworld. \emph{Burke's identification / consubstantiation} supplies the rhetorical bond by which agents become ``substantially one, yet a distinct locus of motives.'' \emph{Rickert's reading of Heideggerian dwelling} supplies the ecological, never-finished way a being is conditioned by---and gathered into---the world it inhabits. RODA treats these not as three lenses used in turn but as \emph{two faces and one ground of a single event}: incorporation is the world-face, consubstantiation the social other-face, and dwelling the prior horizon from which both become possible---all unified as moments of one actualization of desire.

\medskip
\noindent\textbf{Calleja's Player Involvement Model (PIM).} The structural vocabulary RODA absorbs as \emph{designed channels} (Step~1 names a dimension only where the world applies pressure):
\begin{itemize}
  \item \textbf{Kinesthetic}---movement and control of the avatar; the player's motor mediation, ripening from symbolic command through mimetic correspondence to ``symbiotic'' fluency. The substrate of the other five.
  \item \textbf{Spatial}---the internalization of game space as an inhabited \emph{place}: orientation, navigation, the felt sense of being-there. The cornerstone of incorporation.
  \item \textbf{Shared}---involvement with other agents, human or artificial: cohabitation, cooperation, competition---the social field of play.
  \item \textbf{Narrative}---the game's story as both \emph{scripted} (delivered to the player) and \emph{emergent} (the player's own ``alterbiography,'' assembled through play).
  \item \textbf{Affective}---the emotional register: mood and atmosphere, and the emotions a game designs for or precipitates.
  \item \textbf{Ludic}---the rule-system and its goals; the structure of choice-under-repercussion, the ``wager'' of play.
\end{itemize}

\medskip
\noindent\textbf{The incorporation criteria (R1--R7).} How RODA judges whether the world-face has actually been reached. R3--R7 are Calleja's own test for \emph{incorporation} as against mere \emph{involvement}---fail any one and the player is only involved; RODA numbers and extends them.
\begin{itemize}
  \item \textbf{R1---Assimilation.} The gameworld is taken into the player's awareness as their immediate surroundings.
  \item \textbf{R2---Embodiment.} The player is, at the same time, embodied \emph{within} that world through the avatar---the converse direction.
  \item \textbf{R3---Simultaneity.} R1 and R2 held \emph{at once}: assimilation and embodiment together. The strictest gatekeeper---the player is joined to the world yet remains a distinct locus (``joined-and-separate'').
  \item \textbf{R4---Cornerstone.} The spatial-and-kinesthetic base: a body oriented in an inhabited place (front/back, left/right; a future ahead, a past behind), on which the rest rests.
  \item \textbf{R5---Blending.} Several involvement dimensions (typically four or five) co-active and \emph{fused}, not attended to in turn.
  \item \textbf{R6---Temporal apex.} Incorporation as the culmination of a sustained arc, deepening the longer it is held; a brief or early beat reaches it only partially.
  \item \textbf{R7---Intrusion-integration.} Real-world stimuli \emph{blend in} rather than shattering the union---the porosity through which what happens in the game can later cross into the player's real-world being (the channel of carry-over).
\end{itemize}

\section{The ontological frame}
\noindent The frame is general: it holds of any action, in a game or out of one. It is laid down once, here, and then invoked only where it does work.
\begin{itemize}
  \item \textbf{Actualization, in two forms.} Every analyzed episode is the actualization of a potential, and Aristotle distinguishes two kinds. \emph{Energeia} in the \emph{narrow} sense is an actuality whose end is \emph{internal} and co-present---the activity is its own completion, self-sustaining (seeing, thinking, living-well: one sees and \emph{has}-seen at once). \emph{Kinesis} proper is the actuality of a potential \emph{as} potential, whose end lies \emph{outside} it---incomplete and self-canceling (building ceases once the house stands). So ``every episode is a motion'' must be qualified: the chain's links are \textit{kin\=eseis}, its nodes are reached actualities, and the terminal act $A_4$ is itself either a \textit{praxis} (end-internal, an \textit{energeia}) or a \textit{poi\=esis} (end-external, a \textit{kin\=esis}).
  \item \textbf{A change of a substance, indexed to time.} Each episode is a \emph{metabol\=e} undergone by a determinate matter--form complex, classifiable by category: substance (\textit{genesis}/\textit{phthora}), quality (\textit{alloi\=osis}), quantity, or place (\textit{phora}). Two substance-series typically run at once---the player (whose perceptual, cognitive, and dispositional states alter, sedimenting as \textit{hexis}) and the world/avatar (altered by the player's \textit{poi\=esis}).
  \item \textbf{The four causes---two analyses, hinged (not a partition).} The four causes belong to a \emph{single} change, so they cannot be parcelled out between two agents. \emph{Designing} and \emph{playing} are two different actions, each with its own complete set of four causes, joined by Aristotle's producing-art / using-art relation (\emph{Phys.} II.2): the designer is the \emph{producing-art} (``knows the matter''), the player the \emph{using-art} (``knows the form''); and the using-art is \emph{architectonic}---the playing is the for-the-sake-of-which the producing serves.
    \begin{itemize}
      \item \emph{The designer's making} (a \textit{poi\=esis}): material $=$ the medium (engine, assets, systems); formal $=$ the design (rules and the intended-experience-as-form); efficient $=$ the designer; final $=$ that the game be played-and-experienced thus---the player's activity-to-come, \emph{anticipated}.
      \item \emph{The player's playing} (a \textit{praxis}): material $=$ the game-as-given (world, avatar, affordances); formal $=$ the form they know and complete (the avatar-form inhabited, the rules grasped as the shape of good play); efficient $=$ the player's own \textit{orexis}; final $=$ their own \textit{orekton}.
      \item \emph{The hinge}: one artifact, two angles---the designer's \emph{end-product} (the final cause of the \emph{making}) \emph{is} the player's \emph{given means} (the material and formal of the \emph{playing}). The designer therefore cannot be the efficient cause of the player's act---one agent cannot supply another's desire---but only authors the player's \emph{world} and \emph{anticipates} the player's end. ``Why now'' is the designer altering the player's material world to solicit a response; ``which way'' is the player's own efficient-and-final. Alignment of the two ends is a smooth design; the player's actual \textit{orekton} \emph{exceeding} the anticipated one is emergent play; its \emph{collapse} into the designed end, one face of compulsion.
    \end{itemize}
  \item \textbf{The three-factor schema.} Within any motion the efficient and final causes operate as a relay: a \emph{relative unmoved mover} (the \textit{orekton}, the desired good, which moves without being moved), a \emph{moved mover} (\textit{orexis}, desire---moved by the good and moving the body), and the \emph{moved} (the animal, by way of the joint).
  \item \textbf{The seven causes, and the four types.} For voluntary human action the efficient--final complex is individuated by Aristotle's seven causes (\emph{Rhetorica} I.10): the \emph{involuntary} (chance; and necessity---compulsion and nature) and the \emph{voluntary} (habit, the dispositional ground; and the desires). From these the four action-\emph{types} derive---Epithymetic (\textit{epithymia}, appetite), Thymotic (\textit{thymos}, spirit), Bouletic (\textit{boul\=esis}, rational wish), and Prohairetic (\textit{prohairesis}, deliberated choice)---with habit underlying all four.
\end{itemize}

\section{The protocol}
\noindent The episode is read once, through the chain. Each component below is stated first as a piece of method; the indented sub-point gives its correlate in the brawl.

\subsection*{Step 0---Frame}
\noindent The framing step fixes the unit of analysis---its scope, scale, and point of entry---names the object of desire driving the pass, and identifies the prior world the player already inhabits.
\begin{itemize}
  \item \textbf{Scale.} Whether the episode is a discrete beat, a protracted sequence, or a moment of ontological change---which sets how much of the chain and the diachronic register the reading must carry.
    \begin{itemize}\item \inb \emph{discrete} (sub-beats: trigger / escalation / climax / mercy).\end{itemize}
  \item \textbf{Entry point.} Where the pass begins---in sense ($A_0$), in imagination ($A_2$), or in thought ($A_3$).
    \begin{itemize}\item \inb \emph{$A_0$-sense}: Bill's shove enters Arthur's perceptual field.\end{itemize}
  \item \textbf{Register.} Synchronic (one directional pass) or diachronic (how passes accumulate and recondition one another).
    \begin{itemize}\item \inb synchronic, with a light diachronic tail.\end{itemize}
  \item \textbf{Orekton} (the object of desire). What the player wants---their own (individual) or the group's (collective/identificatory); a collective object signals an act the social union will source.
    \begin{itemize}\item \inb \textbf{collective}: ``stand with the gang, back Bill.''\end{itemize}
  \item \textbf{Dwelling-horizon.} The prior enworldment the pass draws on---the world and community the player is \emph{already} dwelling in before the episode begins.
    \begin{itemize}\item \inb Arthur already dwells with the gang (camp, shared history, the frontier world); the collective object presupposes this, which is why he wades in without deliberation.\end{itemize}
\end{itemize}

\subsection*{Step 1---Trace the actualization (the chain $A_0\too A_4$)}
\noindent The heart of the analysis. The episode is walked node by node: at each \emph{motion} ($M_n\too A_{n+1}$) the three-factor schema fires (good \too\ desire \too\ body); at each \emph{node} ($A_n$) an actuality is reached. Throughout, one asks what the world \emph{supplies} and what the player's desire \emph{does}, naming a channel only where the world is exerting pressure.
\begin{itemize}
  \item \textbf{$A_0$---the horizon.} The ontological starting-point: a sensible object and a sense-faculty given together within motion and time---the designed world as an inhabited \textit{Umwelt}, the dwelling-horizon into which the player is already thrown.
    \begin{itemize}\item \inb the saloon opening onto the muddy street and a gathering crowd; the world \emph{scripts} Bill's altercation into Arthur's perception (\emph{Spatial, Shared}).\end{itemize}
  \item \textbf{$M_0\too A_1$---perceptual motion.} The object's active power and the sense-organ's receptivity actualize as one shared activity; a hedonic tone is co-given.
    \begin{itemize}\item \inb the combat field registers, its thrill the tone (\emph{Kinesthetic, Affective}).\end{itemize}
  \item \textbf{$A_1$---completed perception.} The perceptual act plus its dual residue: a formal trace (the \emph{resonant aisth\=ema}) and a hedonic-orectic trace (the \emph{resonant epithymia}); a being-affected that is here preservative.
    \begin{itemize}\item \inb the ``Hold R to block'' prompt is perceived---the designer shaping the player's given world so that it \emph{teaches its own form}, the combat skill taking hold through the player's own doing.\end{itemize}
  \item \textbf{$M_1\too A_2$---phantastic motion.} The residue settles into a determinate, retrievable image; prior dispositions can already bias what figure forms.
    \begin{itemize}\item \inb the scene crystallizes into an image of ``the brawl---and the gang as mine.''\end{itemize}
  \item \textbf{$A_2$---the \emph{phantasma} (the pivot).} The charged image on which all later operations work; the seed of identification is laid here, the other becoming present-as-mattering.
    \begin{itemize}\item \inb ``we can handle these fools''---the gang is present as \emph{ours} (\emph{Shared}).\end{itemize}
  \item \textbf{$M_2\too A_3$---cognitive motion.} The image is met by the player's settled opinions-as-dispositions and passed through the cognitive apparatus, where opinion gatekeeps what is taken as real.
    \begin{itemize}\item \inb Arthur's standing gang-loyalty meets the image of the slight.\end{itemize}
  \item \textbf{$A_3$---cognition, \emph{doxa}, and emotion.} The completed cognitive actuality has three moments. \emph{Cognitive orientation}---the mode in which the image is taken up (intellection, memory, or discursive projection). \emph{Doxa}---the committal taking-as-true that ratifies the image as how things stand. And \emph{emotion}---arguably the crucial moment: the \emph{determinate form desire takes} in response to the world now disclosed-\emph{as-such}. Once \textit{doxa} ratifies the situation as (say) an affront, the bare tone of $A_1$ concretizes into a specific passion (anger, fear, pity, shame)---the shape the desire assumes toward the disclosed world, and the immediate spring of the act.
    \begin{itemize}\item \inb \textit{doxa} ratifies the ``\textbf{WE}'' (``we back our own''); the crowd ratifies a counter-``WE'' (``Finish him, Tommy!''); and the emotion that concretizes is \textbf{indignation} at the slight to the gang---the determinate form Arthur's \textit{thymos} takes, here the very engine of what follows (\emph{Shared, Affective, Ludic}).\end{itemize}
  \item \textbf{$M_3\too A_4$---orektikon motion.} The good as judged-and-felt at $A_3$ moves \textit{orexis}, which moves the animal through the joint; \emph{which} species of desire fires is settled here.
    \begin{itemize}\item \inb \textit{thymos}---the spirited, irreducibly social desire that answers a slight on behalf of the group---fires.\end{itemize}
  \item \textbf{$A_4$---completed action (\emph{praxis}).} The chain's terminus, where the four causes converge; it individuates into an action-type and sediments back into \textit{hexis}.
    \begin{itemize}\item \inb block, strike, break the grapple (``BREAK FREE [F]''), dominate---the encounter scripted to be won (\emph{Ludic, Kinesthetic}).\end{itemize}
  \item \textbf{$A_4\too A_0'$---recursion.} The completed act alters the world and deposits a residue in the agent, reconditioning the next pass.
    \begin{itemize}\item \inb a combat disposition is laid down, the gang-bond deepens, a cruelty/restraint calibration is set (\too\ Step~6).\end{itemize}
\end{itemize}

\subsection*{Step 2---Diagnose the union (\emph{rhetorical incorporation})}
\noindent The central event. At the $A_3\too A_4$ transition the chain reaches a substantial union---``one substrate, two in being''---with two faces and one prior ground. The step assesses both faces and asks the decisive question: does the union merely crown the act, or \emph{source} it?
\begin{itemize}
  \item \textbf{World-face---incorporation.} The player one-with the world they act in: the gameworld assimilated into awareness and the player embodied within it, held at once; gauged by Calleja's criteria R1--R7.
    \begin{itemize}\item \inb R1--R5 met (arena absorbed; Arthur-as-me; joined-yet-distinct; spatial--kinesthetic cornerstone; several dimensions blended); R6 (temporal apex) only \emph{partial} (an early-tutorial beat); R7 (intrusion) not in evidence.\end{itemize}
  \item \textbf{The world-face's ground---dwelling.} What incorporation \emph{is}, ontologically: the player dwelling in the designed world, constituted \emph{with} it rather than merely receiving it---the world an active party that makes claims.
    \begin{itemize}\item \inb the gang-world and frontier are the ground Arthur acts from, not a backdrop he observes.\end{itemize}
  \item \textbf{Recalcitrance---\emph{thing} vs \emph{object}.} Whether the world pushes back and makes claims (a \emph{thing}, and so genuinely \emph{dwelt}) or yields on demand (an \emph{object}, merely \emph{occupied})---the measure of whether a design affords dwelling at all.
    \begin{itemize}\item \inb \textbf{high} at the bodily level (mud, fists, the grapple that resists), which is what makes the brawl feel \emph{embodied} rather than watched---yet \textbf{low} at the level of outcome (the win is foreclosed).\end{itemize}
  \item \textbf{Other-face---consubstantiation.} The player one-with the avatar they act \emph{as} and the collective they act \emph{with}---substantially one, yet a distinct locus of motives.
    \begin{itemize}\item \inb \textbf{genuine} with the gang (acting together makes consubstantial); a \textbf{counter}-union forms around the crowd backing Tommy.\end{itemize}
  \item \textbf{The never-stilled union.} The union is \emph{won, not had}---a perpetually re-achieved accomplishment held against constant slippage, never a fixed possession.
    \begin{itemize}\item \inb the incorporation is still \emph{building} and fragile; it lapses the instant the beat ends.\end{itemize}
  \item \textbf{The decisive question---does the union \emph{source} the act?} Whether the union is merely reached or is what generates the next action.
    \begin{itemize}\item \inb \textbf{yes}---the brawl is sourced by the gang-consubstantiation: Arthur fights because he is one with the gang, not from a private slight.\end{itemize}
\end{itemize}

\subsection*{Step 3---Read the motive-rhetoric (Burke)}
\noindent\emph{Set aside for now.}

\subsection*{Step 4---Cause \& type}
\noindent The causal diagnosis individuates the action within the seven causes: whether it is the agent's own or the world acting on them, and, if theirs, which desire sources it.
\begin{itemize}
  \item \textbf{Cause.} Agent-caused (voluntary, desire-sourced) or not (chance, nature, or \emph{force}); the \emph{force} test asks whether the moving principle is external, ``contrary to the desire or reason of the agent.''
    \begin{itemize}\item \inb \textbf{agent-caused and voluntary}---the \emph{force} test fails (Arthur acts from his own spirit). What chance there is (damage variance) is authored, not true chance.\end{itemize}
  \item \textbf{Action-type.} The predominant source-desire: appetitive, thymotic, bouletic, or prohairetic.
    \begin{itemize}\item \inb \textbf{Thymotic} (a reflex floor below, deliberate timing in the blocks above); the gang-loyalty habit is the ground it fires through, not a competing type.\end{itemize}
  \item \textbf{Agency-locus.} Whether the player's freedom is over \emph{whether}, \emph{how}, or only the \emph{framing} of the act.
    \begin{itemize}\item \inb over the \emph{how}, not the whether---the outcome is foreclosed.\end{itemize}
  \item \textbf{Hexis-mode.} Whether the disposition cultivated is procedural skill (\textit{techn\=e}) or evaluative character (ethical).
    \begin{itemize}\item \inb \textbf{both}---a combat \textit{techn\=e} taught live, and an ethical fork staged at the mercy beat.\end{itemize}
\end{itemize}

\subsection*{Step 5---Phantasmatic \& affective specifics}
\noindent A pass over the image-and-affect machinery, including what does \emph{not} fire---the instruments are diagnostic by their silence too.
\begin{itemize}
  \item \textbf{Phantasia-conscription.} When a world withholds depiction and conscripts the player's own imagination to author the unshown, the content fixed by their settled dispositions.
    \begin{itemize}\item \inb \textbf{minimal}---the brawl is shown, not withheld (the contrast that sets up the capstone, where withholding is the engine).\end{itemize}
  \item \textbf{Pathetic loop.} The way an aroused affect biases the tone, the forming image, and the doxa-gate of subsequent passes.
    \begin{itemize}\item \inb crowd bloodlust and thrill bias toward excess; the mercy appeal is the corrective.\end{itemize}
  \item \textbf{Agent\,$\to$\,patient inversion.} A forcing of the player from the active pole back to the passive pole of being-affected.
    \begin{itemize}\item \inb only a \emph{micro}-instance: the grapple briefly \emph{gathers} Arthur---the world claims him---before he mashes free; a small taste of what the capstone weaponizes.\end{itemize}
\end{itemize}

\subsection*{Step 6---Diachronic \& carry-over}
\noindent What the completed act leaves behind, and---at the largest scale---whether it follows the player out of the game.
\begin{itemize}
  \item \textbf{Hexis sedimented.} The skill and character the act deposits for future passes.
    \begin{itemize}\item \inb a combat \textit{techn\=e}, a deepened gang-bond, and a cruelty/restraint calibration feeding the Honor system.\end{itemize}
  \item \textbf{Stimmung for the next pass.} The mood the residue carries into the next horizon.
    \begin{itemize}\item \inb readiness-for-violence and a camaraderie-glow.\end{itemize}
  \item \textbf{Virtual\,$\to$\,actual carry-over.} Whether the sedimented union supplants the player's real-world dispositional ground---a change in how they \emph{dwell}, the deepest sense in which the work is \emph{rhetorical}.
    \begin{itemize}\item \inb \textbf{not at this scale}: the brawl is a building block of the tutorial's engineered disposition---the very ground the capstone will later supplant.\end{itemize}
\end{itemize}

\subsection*{Step 7---Diagnosis (the synthesis)}
\begin{quote}\small
The saloon brawl is a \textbf{Thymotic act sourced by the consubstantial union with the gang}---ontologically a designer-triggered \textit{alloi\=osis} of the player's combat-\textit{hexis}, drawn from a prior dwelling Arthur already shares with the gang. Two four-cause analyses meet here. The player's playing has its material in the given world (saloon, mud, the avatar Arthur), its form in the melee they enact, its efficient cause in their own spirited desire, and its end in their \textit{orekton} (back the gang). The designer's making supplied that world and form and \emph{anticipated} that end---scripting Bill's shove to occasion the act, foreclosing the win---so the designer's anticipated end and the player's actual one \emph{coincide}: a tutorial built so the player wants exactly what it teaches. At $A_3$ the situation is ratified as an affront and the desire takes determinate form as \emph{indignation}; the ``Hold R to block'' prompt is the world teaching its own form, a \textit{techn\=e} taking hold through the player's own doing; the brawl's bodily recalcitrance (mud, the grapple that resists) makes it \emph{dwelt}, not merely watched, even as its outcome is foreclosed. Its motive is rhetorically rendered as \emph{provoked loyalty}, concealing Arthur's own violence-\textit{hexis}; it cultivates a \textit{techn\=e} while staging an ethical fork (the mercy beat), sedimenting both---the dispositional ground the capstone will later supplant.
\end{quote}

\vspace{4pt}
\noindent\rule{\linewidth}{0.4pt}\\[2pt]
\footnotesize Method: \texttt{CANONICAL-METHODOLOGY-v5.md}. Analysis: \texttt{VALIDATION-v5/v4/v3-barfight.md}. Aristotle / Rickert loci PDF-verified in \texttt{verbatim\_passages.md}. ``RODA'' and ``rhetorical incorporation'' confirmed 2026-06-09; the \emph{actualization of desire} names the generative spine.
\end{document}
